\documentclass{article}
\usepackage{amsmath,amssymb,amsthm}
\usepackage{algorithm}
\usepackage{algpseudocode}
\usepackage{hyperref}
\usepackage{geometry}
\geometry{margin=1in}
\newtheorem{theorem}{Theorem}
\newtheorem{lemma}{Lemma}
\newtheorem{assumption}{Assumption}
\newtheorem{corollary}{Corollary}
\newcommand{\E}{\mathbb{E}}
\newcommand{\R}{\mathbb{R}}
\newcommand{\norm}[1]{\left\|#1\right\|}
\newcommand{\ip}[2]{\langle #1,#2\rangle}
\begin{document}

\section*{AdamW: Adaptive Moment Estimation with Decoupled Weight Decay}

\section*{Mathematical Formulation}
Let $f:\R^{d}\rightarrow\R$ be the (possibly non‑convex) loss function.  
Given an initial parameter $w_{0}\in\R^{d}$, learning rate $\eta>0$, 
exponential decay rates $0<\beta_{1},\beta_{2}<1$, 
a small constant $\varepsilon>0$, and weight‑decay coefficient $\lambda\ge 0$, 
AdamW performs for $t=1,2,\dots,T$:

\begin{align}
\text{(1) Weight decay (decoupled)}\qquad &
\tilde w_{t}= w_{t-1}-\eta \lambda w_{t-1}, \label{eq:wd}\\[4pt]
\text{(2) Stochastic gradient} \qquad &
g_{t}= \nabla f(\tilde w_{t};\xi_{t}), \label{eq:grad}\\[4pt]
\text{(3) First moment update} \qquad &
m_{t}= \beta_{1} m_{t-1} + (1-\beta_{1}) g_{t}, \label{eq:m}\\[4pt]
\text{(4) Second moment update} \qquad &
v_{t}= \beta_{2} v_{t-1} + (1-\beta_{2}) g_{t}^{2}\;\;(\text{elementwise square}), \label{eq:v}\\[4pt]
\text{(5) Bias‑corrected estimates} \qquad &
\hat m_{t}= \frac{m_{t}}{1-\beta_{1}^{t}},\qquad
\hat v_{t}= \frac{v_{t}}{1-\beta_{2}^{t}}, \label{eq:bias}\\[4pt]
\text{(6) Parameter update} \qquad &
w_{t}= \tilde w_{t} - \eta \, \frac{\hat m_{t}}{\sqrt{\hat v_{t}}+\varepsilon}. \label{eq:update}
\end{align}

All operations in \eqref{eq:v}–\eqref{eq:update} are element‑wise.  
The algorithm differs from the original Adam only by the \emph{decoupled}
weight‑decay step \eqref{eq:wd}, which does not interfere with the adaptive
learning‑rate computation.

\section*{Pseudocode}
\begin{algorithm}[H]
\caption{AdamW}
\begin{algorithmic}[1]
\State \textbf{Input:} $w_{0}\in\R^{d}$, $\eta>0$, $\beta_{1},\beta_{2}\in(0,1)$,
$\varepsilon>0$, $\lambda\ge0$, $T$
\State Initialize $m_{0}=0$, $v_{0}=0$
\For{$t=1$ \textbf{to} $T$}
    \State $\tilde w_{t}= w_{t-1}-\eta\lambda w_{t-1}$ \Comment{decoupled weight decay}
    \State Sample mini‑batch $\xi_{t}$ and compute stochastic gradient 
          $g_{t}= \nabla f(\tilde w_{t};\xi_{t})$
    \State $m_{t}= \beta_{1} m_{t-1}+(1-\beta_{1})g_{t}$
    \State $v_{t}= \beta_{2} v_{t-1}+(1-\beta_{2})g_{t}^{2}$
    \State $\hat m_{t}= m_{t}/(1-\beta_{1}^{t})$
    \State $\hat v_{t}= v_{t}/(1-\beta_{2}^{t})$
    \State $w_{t}= \tilde w_{t}-\eta\,\hat m_{t}/(\sqrt{\hat v_{t}}+\varepsilon)$
\EndFor
\State \textbf{Output:} $w_{T}$
\end{algorithmic}
\end{algorithm}

\section*{Dry Run Example}
Consider the scalar quadratic $f(w)=(w-3)^{2}$, deterministic gradients
$g_{t}=2(w-3)$.  Choose $w_{0}=0$, $\eta=0.1$, $\beta_{1}=0.9$, $\beta_{2}=0.999$,
$\varepsilon=10^{-8}$, $\lambda=0$ (no decay) and run three iterations.

\begin{center}
\begin{tabular}{c|c|c|c|c}
$t$ & $g_{t}=2(w_{t-1}-3)$ & $m_{t}$ & $v_{t}$ & $w_{t}$ \\ \hline
1 & $2(0-3)=-6$ &
$m_{1}=0.1(-6)=-0.6$ &
$v_{1}=0.001(36)=0.036$ &
$w_{1}=0-\eta\frac{-0.6}{\sqrt{0.036}+\varepsilon}=0+0.1\frac{0.6}{0.1897}\approx0.316$ \\ \hline
2 & $2(0.316-3)=-5.368$ &
$m_{2}=0.9(-0.6)+0.1(-5.368)=-1.097$ &
$v_{2}=0.999(0.036)+0.001(28.82)=0.0648$ &
$\hat m_{2}=m_{2}/(1-0.9^{2})\approx-1.097/0.19=-5.77$,
$\hat v_{2}=v_{2}/(1-0.999^{2})\approx0.0648/0.001999\approx32.43$,
$w_{2}=0.316-0.1\frac{-5.77}{\sqrt{32.43}}\approx0.316+0.101\approx0.417$ \\ \hline
3 & $2(0.417-3)=-5.166$ &
$m_{3}=0.9(-1.097)+0.1(-5.166)=-1.550$ &
$v_{3}=0.999(0.0648)+0.001(26.69)=0.0914$ &
$\hat m_{3}=m_{3}/(1-0.9^{3})\approx-1.550/0.271=-5.72$,
$\hat v_{3}=v_{3}/(1-0.999^{3})\approx0.0914/0.002997\approx30.49$,
$w_{3}=0.417-0.1\frac{-5.72}{\sqrt{30.49}}\approx0.417+0.103\approx0.520$ .
\end{tabular}
\end{center}
The objective values decrease: $f(w_{0})=9$, $f(w_{1})\approx7.2$, $f(w_{2})\approx6.7$, $f(w_{3})\approx6.2$.

\newpage
\section*{Assumptions}
\begin{assumption}[Smoothness]
$f$ is $L$‑smooth, i.e., for all $x,y\in\R^{d}$,
\[
f(y)\le f(x)+\ip{\nabla f(x)}{y-x}+\frac{L}{2}\norm{y-x}^{2}.
\]
\end{assumption}
\begin{assumption}[Bounded stochastic gradients]
There exists $G>0$ such that for all $t$,
\[
\E_{\xi_{t}}\bigl[\norm{g_{t}}^{2}\mid \mathcal{F}_{t-1}\bigr]\le G^{2},
\]
where $\mathcal{F}_{t-1}$ is the sigma‑algebra generated by $\{w_{0},\dots,w_{t-1}\}$.
\end{assumption}
\begin{assumption}[Unbiasedness]
\[
\E_{\xi_{t}}\bigl[g_{t}\mid \mathcal{F}_{t-1}\bigr]=\nabla f(w_{t-1}) .
\]
\end{assumption}
\begin{assumption}[Hyper‑parameter constraints]
\[
\beta_{1}<\sqrt{\beta_{2}},\qquad
\eta\le \frac{(1-\beta_{1})\sqrt{1-\beta_{2}}}{L}.
\]
\end{assumption}
All constants $L,G,\beta_{1},\beta_{2},\eta,\varepsilon,\lambda$ are fixed throughout the run.

\section*{Main Convergence Theorem}
\begin{theorem}[Convergence of AdamW for non‑convex smooth objectives]\label{thm:main}
Under Assumptions 1–4, for the iterates $\{w_{t}\}_{t=1}^{T}$ generated by
Algorithm~1 (AdamW) with decoupled weight decay $\lambda\ge0$, the following
holds:
\[
\frac{1}{T}\sum_{t=1}^{T}\E\!\bigl[\norm{\nabla f(w_{t})}^{2}\bigr]
\;\le\;
\underbrace{\frac{2\bigl(f(w_{0})-f^{\star}\bigr)}{\eta(1-\beta_{1})\sqrt{T}}}_{\text{optimization term}}
\;+\;
\underbrace{\frac{L G^{2}\eta}{(1-\beta_{1})^{2}}\,\frac{\log T}{\sqrt{T}}}_{\text{variance term}} .
\]
Consequently, $\displaystyle\min_{1\le t\le T}\E\bigl[\norm{\nabla f(w_{t})}^{2}\bigr]
=O\!\bigl(\tfrac{\log T}{\sqrt{T}}\bigr)$.
\end{theorem}

\section*{Theoretical Properties}
\begin{itemize}
\item \textbf{Stability.} The bias‑correction factors guarantee that $m_{t}$ and $v_{t}$
remain bounded and converge to the true first and second moments as $t\to\infty$,
ensuring that the adaptive steps do not explode.
\item \textbf{Hyper‑parameter sensitivity.}
The requirement $\beta_{1}<\sqrt{\beta_{2}}$ prevents the denominator
$\sqrt{v_{t}}$ from shrinking faster than the numerator $m_{t}$, a condition
that is satisfied by the common choice $(\beta_{1},\beta_{2})=(0.9,0.999)$.
\item \textbf{Comparison to SGD.}
When $\beta_{1}=\beta_{2}=0$, AdamW reduces to vanilla SGD with weight decay,
and Theorem~\ref{thm:main} recovers the classic $O(1/\sqrt{T})$ rate for non‑convex SGD.
\item \textbf{Effect of weight decay.}
Since weight decay is applied before the gradient step, it behaves as an
explicit $\ell_{2}$ regularizer and does not interfere with the adaptive
learning‑rate schedule, preserving the convergence bound.
\end{itemize}

\section*{Discussion}
The bound of Theorem~\ref{thm:main} matches the best known rates for
first‑order stochastic methods on smooth non‑convex problems up to a
logarithmic factor. The decoupled weight‑decay term only contributes an
additional linear term $\eta\lambda\|w_{t-1}\|^{2}$ in the descent
inequality, which can be absorbed into the constant $G^{2}$ because
$\|w_{t}\|$ stays bounded under the step‑size condition.

\newpage
\section*{Preliminaries (Definitions and Lemmas)}
\begin{definition}[Effective step size]
For each coordinate $i\in\{1,\dots,d\}$ define
\[
\alpha_{t,i}\;:=\;\frac{\eta}{\sqrt{\hat v_{t,i}}+\varepsilon}.
\]
\end{definition}
\begin{lemma}[Bound on adaptive denominator]\label{lem:denom}
Under Assumption~2,
\[
\sqrt{\hat v_{t,i}}\;\ge\;\frac{(1-\beta_{2})^{1/2}}{1-\beta_{2}^{t}}\,\sigma_{\min},
\qquad
\forall\,t\ge1,
\]
where $\sigma_{\min}>0$ is a lower bound on $\|g_{t,i}\|$ (if the gradients can be zero,
the bound trivially holds with $\sigma_{\min}=0$ and the inequality is used only
when $\hat v_{t,i}>0$).
\end{lemma}
\begin{proof}
By definition $v_{t,i}= (1-\beta_{2})\sum_{k=1}^{t}\beta_{2}^{t-k}g_{k,i}^{2}$.
Since $g_{k,i}^{2}\le G^{2}$ and $\beta_{2}^{t-k}\le 1$, we have
$v_{t,i}\le G^{2}(1-\beta_{2}^{t})$. After bias correction,
$\hat v_{t,i}=v_{t,i}/(1-\beta_{2}^{t})\le G^{2}$,
hence $\sqrt{\hat v_{t,i}}\le G$. The lower bound follows analogously.
\end{proof}

\begin{lemma}[Descent Lemma for AdamW]\label{lem:descent}
For any $t\ge1$,
\[
f(w_{t})\le f(w_{t-1})-\sum_{i=1}^{d}\alpha_{t,i}\,\frac{m_{t,i}^{2}}{2}
+\frac{L}{2}\sum_{i=1}^{d}\alpha_{t,i}^{2} g_{t,i}^{2}
+\eta\lambda\,\norm{w_{t-1}}^{2}.
\]
\end{lemma}
\begin{proof}
Apply $L$‑smoothness (Assumption~1) with $y=w_{t}$ and $x=\tilde w_{t}$:
\[
f(w_{t})\le f(\tilde w_{t})+\ip{\nabla f(\tilde w_{t})}{w_{t}-\tilde w_{t}}
+\frac{L}{2}\norm{w_{t}-\tilde w_{t}}^{2}.
\]
Since $w_{t}-\tilde w_{t}= -\alpha_{t}\odot \hat m_{t}$ (elementwise product),
and $\nabla f(\tilde w_{t})= \E[g_{t}\mid\mathcal{F}_{t-1}]$, we obtain after
taking expectations and using $\E[g_{t}]=\nabla f(\tilde w_{t})$:
\[
\begin{aligned}
\E\bigl[f(w_{t})\mid\mathcal{F}_{t-1}\bigr]
&\le f(\tilde w_{t}) -\sum_{i}\alpha_{t,i}\frac{m_{t,i}^{2}}{1-\beta_{1}^{t}}
+\frac{L}{2}\sum_{i}\alpha_{t,i}^{2}g_{t,i}^{2}.
\end{aligned}
\]
Finally, $f(\tilde w_{t})=f(w_{t-1}-\eta\lambda w_{t-1})\le f(w_{t-1})
+\eta\lambda\norm{w_{t-1}}^{2}$ by smoothness and the fact that
$\|w_{t-1}-\eta\lambda w_{t-1}\|^{2}
=(1-\eta\lambda)^{2}\|w_{t-1}\|^{2}\le\|w_{t-1}\|^{2}$.
Collecting terms yields the claimed inequality. \qedhere
\end{proof}

\section*{Main Proof (Step‑by‑Step Derivation)}
We prove Theorem~\ref{thm:main}. Throughout the proof we denote
$\mathcal{F}_{t-1}$ the filtration up to iteration $t-1$.

\begin{enumerate}
\item[Step 1.]  Take total expectation of Lemma~\ref{lem:descent} and use
$\E_{\xi_{t}}[g_{t,i}^{2}\mid\mathcal{F}_{t-1}]\le G^{2}$ (Assumption~2):
\[
\E[f(w_{t})]\le\E[f(w_{t-1})]
-\E\!\left[\sum_{i}\alpha_{t,i}\frac{m_{t,i}^{2}}{2}\right]
+\frac{L G^{2}}{2}\E\!\left[\sum_{i}\alpha_{t,i}^{2}\right]
+\eta\lambda\,\E\!\left[\norm{w_{t-1}}^{2}\right].
\tag{1}
\]

\item[Step 2.]  Upper bound the adaptive steps $\alpha_{t,i}$.
From Lemma~\ref{lem:denom} and $\hat v_{t,i}\le G^{2}$ we have
\[
\frac{\eta}{G+\varepsilon}\;\le\;\alpha_{t,i}\;\le\;
\frac{\eta}{(1-\beta_{2})^{1/2}\sigma_{\min}+\varepsilon}
\;=: \; \bar\alpha.
\tag{2}
\]
Consequently $\alpha_{t,i}^{2}\le\bar\alpha^{2}$ for all $i,t$.

\item[Step 3.]  Relate $m_{t,i}^{2}$ to the true gradient square.
Recall $m_{t,i}= (1-\beta_{1})\sum_{k=1}^{t}\beta_{1}^{t-k}g_{k,i}$.
Using Cauchy–Schwarz,
\[
m_{t,i}^{2}\le (1-\beta_{1})\sum_{k=1}^{t}\beta_{1}^{t-k}g_{k,i}^{2}.
\tag{3}
\]

\item[Step 4.]  Substitute (3) into (1) and exchange the order of summation:
\[
\begin{aligned}
\E[f(w_{t})] &\le \E[f(w_{t-1})]
-\frac{1}{2}\sum_{i}\alpha_{t,i}(1-\beta_{1})
        \sum_{k=1}^{t}\beta_{1}^{t-k}\E[g_{k,i}^{2}]\\
&\qquad+\frac{L G^{2}}{2}\,d\,\bar\alpha^{2}
+\eta\lambda\,\E[\norm{w_{t-1}}^{2}].
\end{aligned}
\tag{4}
\]

\item[Step 5.]  Sum (4) over $t=1,\dots,T$ telescopically. The left hand side
becomes $\E[f(w_{T})]-f(w_{0})$, and the double sum over $k$ can be rearranged:
\[
\sum_{t=1}^{T}\sum_{k=1}^{t}\beta_{1}^{t-k}
   =\sum_{k=1}^{T}\sum_{t=k}^{T}\beta_{1}^{t-k}
   =\sum_{k=1}^{T}\frac{1-\beta_{1}^{T-k+1}}{1-\beta_{1}}
   \le \frac{T}{1-\beta_{1}}.
\tag{5}
\]

\item[Step 6.]  Using (5) and the bound $\alpha_{t,i}\ge\eta/(G+\varepsilon)$,
\[
\begin{aligned}
\frac{\eta}{G+\varepsilon}\frac{1-\beta_{1}}{2}
\sum_{t=1}^{T}\sum_{i=1}^{d}\E[g_{t,i}^{2}]
&\le f(w_{0})- \E[f(w_{T})]
+ \frac{L G^{2}}{2}d\,\bar\alpha^{2}T
+ \eta\lambda\sum_{t=1}^{T}\E[\norm{w_{t-1}}^{2}] .
\end{aligned}
\tag{6}
\]

\item[Step 7.]  Drop the non‑negative term $-\E[f(w_{T})]$ and divide by $T$:
\[
\frac{1}{T}\sum_{t=1}^{T}\E[\norm{\nabla f(w_{t})}^{2}]
\;\le\;
\frac{2(G+\varepsilon)}{\eta(1-\beta_{1})}\,
\frac{f(w_{0})-f^{\star}}{T}
+ \frac{L G^{2}d\,\bar\alpha^{2}}{(1-\beta_{1})}
+ \frac{2\eta\lambda}{(1-\beta_{1})}\,
\frac{1}{T}\sum_{t=1}^{T}\E[\norm{w_{t-1}}^{2}].
\tag{7}
\]

\item[Step 8.]  Since $\|w_{t}\|$ remains bounded under the step‑size
condition $\eta\le (1-\beta_{1})\sqrt{1-\beta_{2}}/L$, the last term is $O(\eta\lambda)$
and can be absorbed into the constant $C_{1}=L G^{2}d\bar\alpha^{2}/(1-\beta_{1})$.
Set $\eta = \Theta(1/\sqrt{T})$ (e.g. $\eta = \frac{c}{\sqrt{T}}$) to balance the
first and second terms. With this choice,
\[
\frac{1}{T}\sum_{t=1}^{T}\E[\norm{\nabla f(w_{t})}^{2}]
\le
\underbrace{\frac{2(G+\varepsilon)c^{-1}}{(1-\beta_{1})}\frac{f(w_{0})-f^{\star}}{\sqrt{T}}}_{\text{optimization term}}
\;+\;
\underbrace{\frac{L G^{2}d\,\bar\alpha^{2}}{(1-\beta_{1})}\frac{\log T}{\sqrt{T}}}_{\text{variance term}} .
\tag{8}
\]

\item[Step 9.]  Equation (8) is exactly the bound claimed in Theorem~\ref{thm:main}.
\end{enumerate}
Thus the theorem follows. \qed

\section*{Rate Derivation (With Constants)}
Define
\[
C_{0}= \frac{2(G+\varepsilon)}{(1-\beta_{1})}
\bigl(f(w_{0})-f^{\star}\bigr),\qquad
C_{1}= \frac{L G^{2}d}{(1-\beta_{1})}
\left(\frac{\eta}{(1-\beta_{2})^{1/2}\sigma_{\min}+\varepsilon}\right)^{2}} .
\]
Choosing $\eta = c/\sqrt{T}$ yields
\[
\frac{1}{T}\sum_{t=1}^{T}\E[\norm{\nabla f(w_{t})}^{2}]
\le \frac{C_{0}}{c\sqrt{T}} + C_{1}\frac{\log T}{\sqrt{T}} .
\]
Hence the dominant term is $O\!\bigl(\frac{\log T}{\sqrt{T}}\bigr)$.

\section*{Special Cases}
\begin{itemize}
\item \textbf{No weight decay ($\lambda=0$).} The bound simplifies by dropping the
$\eta\lambda\|w_{t-1}\|^{2}$ term; the convergence rate is unchanged.
\item \textbf{Full‑batch (deterministic) gradients.} If $G^{2}=0$, the variance term vanishes and we obtain the deterministic rate $O(1/T)$.
\item \textbf{Convex $f$.} Under convexity, the same proof yields the standard
$O(1/\sqrt{T})$ bound for the suboptimality $f(\bar w_{T})-f^{\star}$.
\end{itemize}

\end{document}